% Options for packages loaded elsewhere
\PassOptionsToPackage{unicode}{hyperref}
\PassOptionsToPackage{hyphens}{url}
\documentclass[
  12pt,
]{report}
\usepackage{xcolor}
\usepackage[margin=1in]{geometry}
\usepackage{amsmath,amssymb}
\setcounter{secnumdepth}{5}
\usepackage{iftex}
\ifPDFTeX
  \usepackage[T1]{fontenc}
  \usepackage[utf8]{inputenc}
  \usepackage{textcomp} % provide euro and other symbols
\else % if luatex or xetex
  \usepackage{unicode-math} % this also loads fontspec
  \defaultfontfeatures{Scale=MatchLowercase}
  \defaultfontfeatures[\rmfamily]{Ligatures=TeX,Scale=1}
\fi
\usepackage{lmodern}
\ifPDFTeX\else
  % xetex/luatex font selection
\fi
% Use upquote if available, for straight quotes in verbatim environments
\IfFileExists{upquote.sty}{\usepackage{upquote}}{}
\IfFileExists{microtype.sty}{% use microtype if available
  \usepackage[]{microtype}
  \UseMicrotypeSet[protrusion]{basicmath} % disable protrusion for tt fonts
}{}
\makeatletter
\@ifundefined{KOMAClassName}{% if non-KOMA class
  \IfFileExists{parskip.sty}{%
    \usepackage{parskip}
  }{% else
    \setlength{\parindent}{0pt}
    \setlength{\parskip}{6pt plus 2pt minus 1pt}}
}{% if KOMA class
  \KOMAoptions{parskip=half}}
\makeatother
\usepackage{color}
\usepackage{fancyvrb}
\newcommand{\VerbBar}{|}
\newcommand{\VERB}{\Verb[commandchars=\\\{\}]}
\DefineVerbatimEnvironment{Highlighting}{Verbatim}{commandchars=\\\{\}}
% Add ',fontsize=\small' for more characters per line
\newenvironment{Shaded}{}{}
\newcommand{\AlertTok}[1]{\textcolor[rgb]{1.00,0.00,0.00}{\textbf{#1}}}
\newcommand{\AnnotationTok}[1]{\textcolor[rgb]{0.38,0.63,0.69}{\textbf{\textit{#1}}}}
\newcommand{\AttributeTok}[1]{\textcolor[rgb]{0.49,0.56,0.16}{#1}}
\newcommand{\BaseNTok}[1]{\textcolor[rgb]{0.25,0.63,0.44}{#1}}
\newcommand{\BuiltInTok}[1]{\textcolor[rgb]{0.00,0.50,0.00}{#1}}
\newcommand{\CharTok}[1]{\textcolor[rgb]{0.25,0.44,0.63}{#1}}
\newcommand{\CommentTok}[1]{\textcolor[rgb]{0.38,0.63,0.69}{\textit{#1}}}
\newcommand{\CommentVarTok}[1]{\textcolor[rgb]{0.38,0.63,0.69}{\textbf{\textit{#1}}}}
\newcommand{\ConstantTok}[1]{\textcolor[rgb]{0.53,0.00,0.00}{#1}}
\newcommand{\ControlFlowTok}[1]{\textcolor[rgb]{0.00,0.44,0.13}{\textbf{#1}}}
\newcommand{\DataTypeTok}[1]{\textcolor[rgb]{0.56,0.13,0.00}{#1}}
\newcommand{\DecValTok}[1]{\textcolor[rgb]{0.25,0.63,0.44}{#1}}
\newcommand{\DocumentationTok}[1]{\textcolor[rgb]{0.73,0.13,0.13}{\textit{#1}}}
\newcommand{\ErrorTok}[1]{\textcolor[rgb]{1.00,0.00,0.00}{\textbf{#1}}}
\newcommand{\ExtensionTok}[1]{#1}
\newcommand{\FloatTok}[1]{\textcolor[rgb]{0.25,0.63,0.44}{#1}}
\newcommand{\FunctionTok}[1]{\textcolor[rgb]{0.02,0.16,0.49}{#1}}
\newcommand{\ImportTok}[1]{\textcolor[rgb]{0.00,0.50,0.00}{\textbf{#1}}}
\newcommand{\InformationTok}[1]{\textcolor[rgb]{0.38,0.63,0.69}{\textbf{\textit{#1}}}}
\newcommand{\KeywordTok}[1]{\textcolor[rgb]{0.00,0.44,0.13}{\textbf{#1}}}
\newcommand{\NormalTok}[1]{#1}
\newcommand{\OperatorTok}[1]{\textcolor[rgb]{0.40,0.40,0.40}{#1}}
\newcommand{\OtherTok}[1]{\textcolor[rgb]{0.00,0.44,0.13}{#1}}
\newcommand{\PreprocessorTok}[1]{\textcolor[rgb]{0.74,0.48,0.00}{#1}}
\newcommand{\RegionMarkerTok}[1]{#1}
\newcommand{\SpecialCharTok}[1]{\textcolor[rgb]{0.25,0.44,0.63}{#1}}
\newcommand{\SpecialStringTok}[1]{\textcolor[rgb]{0.73,0.40,0.53}{#1}}
\newcommand{\StringTok}[1]{\textcolor[rgb]{0.25,0.44,0.63}{#1}}
\newcommand{\VariableTok}[1]{\textcolor[rgb]{0.10,0.09,0.49}{#1}}
\newcommand{\VerbatimStringTok}[1]{\textcolor[rgb]{0.25,0.44,0.63}{#1}}
\newcommand{\WarningTok}[1]{\textcolor[rgb]{0.38,0.63,0.69}{\textbf{\textit{#1}}}}
\setlength{\emergencystretch}{3em} % prevent overfull lines
\providecommand{\tightlist}{%
  \setlength{\itemsep}{0pt}\setlength{\parskip}{0pt}}
\usepackage{bookmark}
\IfFileExists{xurl.sty}{\usepackage{xurl}}{} % add URL line breaks if available
\urlstyle{same}
\hypersetup{
  pdftitle={Ataree Workflows: A Comprehensive eBPF-based Linux Runtime Security Framework},
  pdfauthor={Academic Project Report},
  hidelinks,
  pdfcreator={LaTeX via pandoc}}

\title{Ataree Workflows: A Comprehensive eBPF-based Linux Runtime
Security Framework}
\author{Academic Project Report}
\date{\today}

\begin{document}
\maketitle
\begin{abstract}
As cloud-native environments and microservice architectures continue to
dominate the modern digital landscape, the requirement for robust,
low-overhead introspection within the Linux kernel has grown
exponentially. Standard user-space protections are increasingly evaded
by sophisticated kernel-level exploit methodologies and
hardware-adjacent vulnerabilities limit the efficacy of generic endpoint
security solutions. This report details the theoretical formulation,
architectural design, and deployment implementation of \textbf{Ataree},
a modernized Linux Runtime Security mechanism. Ataree strategically
consolidates the structural documentation robustness of the
\texttt{ataree} project and the low-level execution power of the
\texttt{KShield} agent-manager paradigm. The resulting framework
provides unparalleled introspection via enhanced Berkeley Packet Filter
(eBPF) technologies, leveraging direct hooks into kernel tracepoints,
network stacks, and Linux Security Modules (LSM) to deliver real-time
telemetry streaming and automated enforcement configurations. The system
boasts an extensible TOML-based 0-day detection suite and a deeply
optimized, glassmorphic UX backend dashboard deployed over a secure,
authenticated REST API.
\end{abstract}

{
\setcounter{tocdepth}{2}
\tableofcontents
}
\newpage

\chapter{Introduction \& Core
Overview}\label{introduction-core-overview}

\begin{figure}[h]
\centering
\vspace{5cm}
\caption{Illustration mapping the logical boundary of the Introduction & Core Overview implementation.}
\end{figure}

\textbf{Ataree} is a Linux runtime security project built using eBPF,
combining the robustness of the KShield telemetry manager and agent with
modern documentation and UI methodologies.

It observes kernel activity in real time and streams structured events
securely to an IPS Control Plane (Manager) for inspection, filtering,
analysis, and active enforcement via LSM hooks.

\section{Overview}\label{overview}

At its core, Ataree uses:

\begin{itemize}
\tightlist
\item
  \textbf{eBPF Agent (Rust)}: BPF programs tracing
  \texttt{sys\_enter\_*}, \texttt{esp\_input}, and other core kernel
  tracepoints, effectively mapping raw kernel state to telemetry
  structures.
\item
  \textbf{Python Manager}: A fast backend server ingesting streams of
  telemetry over a REST API and correlating them with runtime
  \texttt{.toml} detectors.
\item
  \textbf{LSM Enforcement}: Actively blocking unauthorized file writes,
  network hooks, and \texttt{ptrace} manipulations without incurring
  context-switching overhead.
\item
  \textbf{Dashboard}: A gorgeous, glassmorphic frontend UI pulling
  metrics securely via REST.
\end{itemize}

\section{Core Features}\label{core-features}

\begin{itemize}
\tightlist
\item
  \textbf{Novel 0-Day Detection}: Accurately intercepting explicit
  mechanisms such as CVE-2026-43284 IPsSec ESP UDP fast-paths and
  CVE-2026-46333 \texttt{pidfd\_getfd} exploit attempts.
\item
  \textbf{Extensible Detection}: Add new logic via
  \texttt{manager/detectors/*.toml} files that are automatically
  synchronized to the kernel and manager state; zero compilation
  required.
\item
  \textbf{Beautiful UX}: A fully modern, beautifully orchestrated,
  natively responsive dark mode display for Security Operators.
\end{itemize}

!!! tip ``Quick Start'' Check out the
\textbf{\href{api-reference.md}{API Reference}} to script against the
manager, or read the \textbf{\href{extensibility.md}{Extensibility
Guide}} to add your very own kernel rules!

\clearpage

\chapter{System Architecture}\label{system-architecture}

\begin{figure}[h]
\centering
\vspace{5cm}
\caption{Illustration mapping the logical boundary of the System Architecture implementation.}
\end{figure}

\section{Manager}\label{manager}

\begin{itemize}
\tightlist
\item
  Runtime: Python \texttt{http.server} (\texttt{manager/app.py})
\item
  Storage: SQLite (\texttt{manager/data/siem.db})
\item
  Detector loading: TOML definitions from \texttt{manager/detectors/}
\item
  Detector authoring guide: \url{detectors.md}
\item
  Dashboard: static HTML/CSS/JS in \texttt{manager/static}
\item
  APIs:

  \begin{itemize}
  \tightlist
  \item
    \texttt{POST\ /api/v1/agents/register}
  \item
    \texttt{GET\ /api/v1/agents}
  \item
    \texttt{POST\ /api/v1/events}
  \item
    \texttt{GET\ /api/v1/events}
  \item
    \texttt{GET\ /api/v1/detections}
  \item
    \texttt{GET\ /api/v1/detectors}
  \item
    \texttt{GET\ /api/v1/metrics/summary}
  \item
    \texttt{GET\ /api/v1/policy}
  \item
    \texttt{GET\ /api/v1/policy/\textless{}agent\_id\textgreater{}}
  \item
    \texttt{POST\ /api/v1/policy/blocked-ipv4}
  \item
    \texttt{DELETE\ /api/v1/policy/blocked-ipv4}
  \item
    \texttt{POST\ /api/v1/policy/blocked-bind-port}
  \item
    \texttt{DELETE\ /api/v1/policy/blocked-bind-port}
  \end{itemize}
\end{itemize}

\section{Agent}\label{agent}

\begin{itemize}
\tightlist
\item
  Runtime: Rust (\texttt{agent/src/main.rs})
\item
  Kernel programs:

  \begin{itemize}
  \tightlist
  \item
    \texttt{kprobe/\_\_x64\_sys\_execve} -\textgreater{} execution
    telemetry
  \item
    \texttt{kprobe/\_\_x64\_sys\_openat} -\textgreater{} file access
    telemetry
  \item
    \texttt{kprobe/\_\_x64\_sys\_openat2} -\textgreater{} file access
    telemetry when available
  \item
    \texttt{kprobe/\_\_x64\_sys\_ptrace} -\textgreater{} ptrace
    telemetry
  \item
    \texttt{kprobe/\_\_x64\_sys\_init\_module} -\textgreater{} module
    load telemetry when available
  \item
    \texttt{kprobe/\_\_x64\_sys\_finit\_module} -\textgreater{} module
    load telemetry when available
  \item
    \texttt{lsm/socket\_connect} -\textgreater{} inline IPv4 destination
    blocking
  \item
    \texttt{lsm/socket\_bind} -\textgreater{} inline local bind-port
    blocking
  \item
    \texttt{lsm/ptrace\_access\_check} -\textgreater{} inline ptrace
    deny policy
  \item
    \texttt{lsm/bprm\_check\_security} -\textgreater{} inline selected
    exec-path blocking
  \item
    \texttt{lsm/file\_permission} -\textgreater{} inline protected
    write-target blocking
  \end{itemize}
\item
  Maps:

  \begin{itemize}
  \tightlist
  \item
    \texttt{EVENTS} (ringbuf): kernel -\textgreater{} userspace events
  \item
    \texttt{BLOCKED\_IPV4} (hash): userspace policy -\textgreater{}
    \texttt{socket\_connect} enforcement
  \item
    \texttt{BLOCKED\_BIND\_PORTS} (hash): userspace policy
    -\textgreater{} \texttt{socket\_bind} enforcement
  \item
    \texttt{BLOCKED\_EXEC\_PATHS} (hash): userspace policy
    -\textgreater{} \texttt{bprm\_check\_security} enforcement
  \item
    \texttt{PROTECTED\_WRITE\_TARGETS} (hash): userspace policy
    -\textgreater{} \texttt{file\_permission} enforcement
  \item
    \texttt{DENY\_PTRACE} (hash): userspace policy -\textgreater{}
    \texttt{ptrace\_access\_check} enforcement
  \end{itemize}
\end{itemize}

\section{Event Flow}\label{event-flow}

\begin{enumerate}
\def\labelenumi{\arabic{enumi}.}
\tightlist
\item
  Kernel eBPF programs emit raw events to the ring buffer.
\item
  Agent reads ring buffer events and converts them to JSON.
\item
  Agent POSTs raw events to \texttt{/api/v1/events}.
\item
  Manager stores raw events in the \texttt{events} table.
\item
  Manager evaluates each event against loaded detector definitions.
\item
  Matching detectors produce derived detection hits in the
  \texttt{detections} table.
\item
  Dashboard renders both raw events and triggered detections.
\end{enumerate}

\section{Policy Flow}\label{policy-flow}

\begin{enumerate}
\def\labelenumi{\arabic{enumi}.}
\tightlist
\item
  Operator creates or updates policy in the dashboard.
\item
  Manager persists policy in SQLite.
\item
  Agent polls
  \texttt{/api/v1/policy/\textless{}agent\_id\textgreater{}}.
\item
  Agent rewrites the relevant eBPF policy maps.
\item
  LSM hooks return \texttt{-EPERM} for blocked actions.
\item
  The same LSM hooks still emit ring-buffer events so policy hits are
  visible in the UI and can trigger detectors.
\end{enumerate}

\section{Detector Flow}\label{detector-flow}

\begin{enumerate}
\def\labelenumi{\arabic{enumi}.}
\tightlist
\item
  Operator or developer adds a TOML detector under
  \texttt{manager/detectors/}.
\item
  Manager reloads detector definitions when the directory contents
  change.
\item
  Incoming events are evaluated against detector match criteria such as:

  \begin{itemize}
  \tightlist
  \item
    \texttt{event\_types}
  \item
    \texttt{actions}
  \item
    \texttt{comm\_in}
  \item
    \texttt{subject\_prefixes}
  \item
    \texttt{subject\_contains}
  \item
    \texttt{dst\_ports}
  \item
    \texttt{dst\_ips}
  \item
    \texttt{arg0\_in}
  \item
    \texttt{arg1\_in}
  \end{itemize}
\item
  Matching events create persisted detection hits.
\item
  Dashboard shows detector hit volume, hit details, and per-detector hit
  counts.
\end{enumerate}

\clearpage

\chapter{Security Dashboard \& Controller API
Reference}\label{security-dashboard-controller-api-reference}

\begin{figure}[h]
\centering
\vspace{5cm}
\caption{Illustration mapping the logical boundary of the Security Dashboard & Controller API Reference implementation.}
\end{figure}

The Ataree backend manager exposes a RESTful API over HTTP to interact
with agents, fetch metrics, and manage LSM enforcement policies. Below
is the API reference.

\begin{quote}
\textbf{Note on Authentication:} Most endpoints used by the dashboard UI
require HTTP Basic Authentication (\texttt{admin:secret}), with the
exception of \texttt{/healthz}, \texttt{/api/v1/agents/register}, and
\texttt{/api/v1/policy/\{agent\_id\}}.
\end{quote}

\begin{center}\rule{0.5\linewidth}{0.5pt}\end{center}

\section{Agent \& Event Ingestion}\label{agent-event-ingestion}

\subsection{\texorpdfstring{\texttt{POST\ /api/v1/agents/register}}{POST /api/v1/agents/register}}\label{post-apiv1agentsregister}

Registers a new eBPF agent. - \textbf{Payload:}
\texttt{\{"hostname":\ "...",\ "ip":\ "...",\ "kernel":\ "...",\ "version":\ "..."\}}
- \textbf{Returns:}
\texttt{\{"agent\_id":\ "\textless{}uuid\textgreater{}"\}}

\subsection{\texorpdfstring{\texttt{POST\ /api/v1/events}}{POST /api/v1/events}}\label{post-apiv1events}

Submit a batch of raw telemetry events gathered by the eBPF agent from
the kernel ring buffer. - \textbf{Payload:}
\texttt{\{"events":\ {[}\ \{\ ...\ \}\ {]}\}} - \textbf{Returns:}
\texttt{\{"written":\ int\}}

\begin{center}\rule{0.5\linewidth}{0.5pt}\end{center}

\section{Dashboard Data}\label{dashboard-data}

\subsection{\texorpdfstring{\texttt{GET\ /api/v1/metrics/summary}}{GET /api/v1/metrics/summary}}\label{get-apiv1metricssummary}

Returns overarching aggregation statistics for the dashboard. -
\textbf{Returns:} Total occurrences, blocked events, severity breakdown,
and loaded policy counts.

\subsection{\texorpdfstring{\texttt{GET\ /api/v1/agents}}{GET /api/v1/agents}}\label{get-apiv1agents}

Returns a list of all known registered agents and their last seen
statuses.

\subsection{\texorpdfstring{\texttt{GET\ /api/v1/events}}{GET /api/v1/events}}\label{get-apiv1events}

Returns a paginated list of recent raw telemetry events. - \textbf{Query
Params:} \texttt{limit} (default: 200, max: 2000)

\subsection{\texorpdfstring{\texttt{GET\ /api/v1/detections}}{GET /api/v1/detections}}\label{get-apiv1detections}

Returns a list of events that matched active TOML detectors (Triggered
Detectors). - \textbf{Query Params:} \texttt{limit} (default: 200, max:
2000)

\subsection{\texorpdfstring{\texttt{GET\ /api/v1/detectors}}{GET /api/v1/detectors}}\label{get-apiv1detectors}

Returns all active loaded rules from the
\texttt{manager/detectors/*.toml} directory, with their current live hit
counts.

\begin{center}\rule{0.5\linewidth}{0.5pt}\end{center}

\section{Policy \& LSM Enforcement}\label{policy-lsm-enforcement}

Policies dictating what gets blocked (IPv4s, bind ports, exec paths,
file writes, ptrace) can be dynamically configured.

\subsection{\texorpdfstring{\texttt{GET\ /api/v1/policy}}{GET /api/v1/policy}}\label{get-apiv1policy}

Returns all globally defined policies.

\subsection{\texorpdfstring{\texttt{GET\ /api/v1/policy/\{agent\_id\}}}{GET /api/v1/policy/\{agent\_id\}}}\label{get-apiv1policyagent_id}

Returns all active policies tailored to a specific agent's scope. The
agents constantly poll this endpoint.

\subsection{Modifying Policies}\label{modifying-policies}

You can add or remove policies via \texttt{POST} or \texttt{DELETE}
across various sub-routes: - \texttt{/api/v1/policy/blocked-ipv4} -
\texttt{/api/v1/policy/blocked-bind-port} -
\texttt{/api/v1/policy/blocked-exec-path} -
\texttt{/api/v1/policy/protected-write-target} -
\texttt{/api/v1/policy/deny-ptrace}

\textbf{POST Example (Block an IP):}

\begin{Shaded}
\begin{Highlighting}[]
\FunctionTok{\{}
  \DataTypeTok{"ip"}\FunctionTok{:} \StringTok{"203.0.113.10"}\FunctionTok{,}
  \DataTypeTok{"enabled"}\FunctionTok{:} \KeywordTok{true}\FunctionTok{,}
  \DataTypeTok{"agent\_id"}\FunctionTok{:} \StringTok{""} 
\FunctionTok{\}}
\end{Highlighting}
\end{Shaded}

\textbf{DELETE Example:}

\begin{Shaded}
\begin{Highlighting}[]
\FunctionTok{\{}
  \DataTypeTok{"ip"}\FunctionTok{:} \StringTok{"203.0.113.10"}\FunctionTok{,}
  \DataTypeTok{"agent\_id"}\FunctionTok{:} \StringTok{""} 
\FunctionTok{\}}
\end{Highlighting}
\end{Shaded}

\emph{(Passing an empty \texttt{agent\_id} scopes the rule globally.)}

\clearpage

\chapter{Out-of-the-Box Protection: Preloaded
Detectors}\label{out-of-the-box-protection-preloaded-detectors}

\begin{figure}[h]
\centering
\vspace{5cm}
\caption{Illustration mapping the logical boundary of the Out-of-the-Box Protection: Preloaded Detectors implementation.}
\end{figure}

Here are all the out-of-the-box detectors available in Ataree:

\begin{itemize}
\tightlist
\item
  \textbf{Blocked Bind Port}: Highlights bind attempts denied by port
  policy.
\item
  \textbf{Blocked Egress Connection}: Highlights outbound network
  connections blocked by policy.
\item
  \textbf{Blocked Exec Path}: Highlights execution denied by selected
  exec path policy.
\item
  \textbf{Blocked Ptrace}: Highlights ptrace access denied by policy.
\item
  \textbf{Blocked Sensitive Write}: Highlights denied writes to
  protected targets.
\item
  \textbf{BPF Syscall Use}: Flags direct bpf syscall use from userspace.
\item
  \textbf{CVE-2026-43284 ESP-in-UDP Splice}: Flags ESP packets parsed by
  esp\_input, catching potential exploitation of the SKBFL\_SHARED\_FRAG
  decrypt-in-place bug.
\item
  \textbf{CVE-2026-46333 Exploitation Attempt}: Flags usage of
  pidfd\_getfd to copy file descriptors from privileged processes,
  indicating the ssh-keysign-pwn race condition.
\item
  \textbf{Interactive Shell Execution}: Flags execution of common
  interactive shells.
\item
  \textbf{Kernel Module Load Attempt}: Flags runtime kernel module
  loading attempts.
\item
  \textbf{Mount Activity}: Flags mount syscall activity for manual
  review.
\item
  \textbf{Kernel Memory Recon}: Flags reads of /proc/kcore which often
  indicate kernel memory reconnaissance.
\item
  \textbf{Ptrace Attach Attempt}: Flags ptrace activity that can
  indicate live process tampering or debugging.
\item
  \textbf{Repeated Bind Denials}: Flags repeated blocked bind attempts
  against the same port and process.
\item
  \textbf{Sensitive Shadow Access}: Flags access attempts against
  /etc/shadow.
\item
  \textbf{Setuid To Root}: Flags attempts to change effective uid to 0.
\end{itemize}

\clearpage

\chapter{Extensibility Guide: Writing Custom
Rules}\label{extensibility-guide-writing-custom-rules}

\begin{figure}[h]
\centering
\vspace{5cm}
\caption{Illustration mapping the logical boundary of the Extensibility Guide: Writing Custom Rules implementation.}
\end{figure}

Ataree is designed from the ground up to be fully extensible with
minimal friction. This guide illustrates how you can extend the platform
to detect novel threats.

\section{1. Zero-Compile Exploit
Detection}\label{zero-compile-exploit-detection}

Unlike traditional kernel models that require complex re-compilation
targeting specific kernel headers to identify new exploits,
\textbf{Ataree} allows rapid, declarative \texttt{.toml} syntax
matchers.

When the agent intercepts activity (like
\texttt{sys\_enter\_pidfd\_getfd} or \texttt{esp\_input}), it streams it
immediately to the backend Manager pipeline.

\subsection{The Detector Pipeline}\label{the-detector-pipeline}

\begin{Shaded}
\begin{Highlighting}[]
\NormalTok{graph TD}
\NormalTok{    A[Attack/Exploit Trigger] {-}{-}\textgreater{} B(eBPF Tracepoint)}
\NormalTok{    B {-}{-}\textgreater{} C(BPF Ring Buffer)}
\NormalTok{    C {-}{-}\textgreater{} D[Ataree Rust Agent]}
\NormalTok{    D {-}{-} REST POST {-}{-}\textgreater{} E[Ataree Python Manager]}
\NormalTok{    E {-}{-}\textgreater{} F\{TOML Detectors Engine\}}
\NormalTok{    F {-}{-} Match {-}{-}\textgreater{} G[Triggered Detections Table]}
\NormalTok{    F {-}{-} Clean {-}{-}\textgreater{} H[Normal Telemetry Stream]}
\end{Highlighting}
\end{Shaded}

\section{2. Authoring a New Detector}\label{authoring-a-new-detector}

The Ataree manager dynamically loads \texttt{.toml} detector files from
the \texttt{manager/detectors/} directory. 2. Map it to a human-readable
string inside \texttt{convert\_event()} (e.g.,
\texttt{"my\_new\_event"}). 3. Rebuild the agent using
\texttt{bash\ scripts/build\_agent.sh}.

\begin{center}\rule{0.5\linewidth}{0.5pt}\end{center}

The Unified Shield manager dynamically loads \texttt{.toml} detector
files from the \texttt{manager/detectors/} directory.

\subsection{Detector Structure}\label{detector-structure}

You do not need to modify any code to add a detector. Simply create a
file such as \texttt{my-detector.toml}:

\begin{Shaded}
\begin{Highlighting}[]
\DataTypeTok{id} \OperatorTok{=} \StringTok{"my{-}new{-}rule"}
\DataTypeTok{name} \OperatorTok{=} \StringTok{"My Novel Threat Detector"}
\DataTypeTok{description} \OperatorTok{=} \StringTok{"Flags this specific telemetry signature."}
\DataTypeTok{severity} \OperatorTok{=} \StringTok{"critical"}
\DataTypeTok{tags} \OperatorTok{=} \OperatorTok{[}\StringTok{"custom"}\OperatorTok{,} \StringTok{"threat"}\OperatorTok{]}
\DataTypeTok{summary\_template} \OperatorTok{=} \StringTok{"\{comm\} triggered the custom rule on \{subject\}"}

\KeywordTok{[match]}
\DataTypeTok{event\_types} \OperatorTok{=} \OperatorTok{[}\StringTok{"my\_new\_event"}\OperatorTok{]}
\DataTypeTok{comm\_in} \OperatorTok{=} \OperatorTok{[}\StringTok{"malware"}\OperatorTok{,} \StringTok{"evil"}\OperatorTok{]}
\end{Highlighting}
\end{Shaded}

\subsection{Real-Time Loading}\label{real-time-loading}

The manager automatically synchronizes and reloads \texttt{.toml} files
upon modification without downtime. Just save the file inside
\texttt{manager/detectors} and watch the Dashboard for hits!

\clearpage

\chapter{Detector Engine Capabilities \&
Syntax}\label{detector-engine-capabilities-syntax}

\begin{figure}[h]
\centering
\vspace{5cm}
\caption{Illustration mapping the logical boundary of the Detector Engine Capabilities & Syntax implementation.}
\end{figure}

KShield detectors are TOML files in \texttt{manager/detectors/}. They
match normalized events after the manager stores raw agent events.
Adding a detector does not require rebuilding the agent.

\section{Quick Start}\label{quick-start}

\begin{enumerate}
\def\labelenumi{\arabic{enumi}.}
\tightlist
\item
  Create a new \texttt{*.toml} file under \texttt{manager/detectors/}.
\item
  Give it a unique \texttt{id}.
\item
  Choose match fields under \texttt{{[}match{]}}.
\item
  Start or refresh the manager. The manager reloads detectors when files
  change.
\item
  Generate matching telemetry and check \texttt{/api/v1/detections} or
  the dashboard.
\end{enumerate}

Example:

\begin{Shaded}
\begin{Highlighting}[]
\DataTypeTok{id} \OperatorTok{=} \StringTok{"suspicious{-}shell{-}from{-}temp"}
\DataTypeTok{name} \OperatorTok{=} \StringTok{"Suspicious Shell From Temp"}
\DataTypeTok{description} \OperatorTok{=} \StringTok{"Flags shell execution where the executed path is under /tmp."}
\DataTypeTok{severity} \OperatorTok{=} \StringTok{"warning"}
\DataTypeTok{tags} \OperatorTok{=} \OperatorTok{[}\StringTok{"execution"}\OperatorTok{,} \StringTok{"suspicious{-}path"}\OperatorTok{]}
\DataTypeTok{summary\_template} \OperatorTok{=} \StringTok{"\{comm\} executed \{subject\}"}

\KeywordTok{[match]}
\DataTypeTok{event\_types} \OperatorTok{=} \OperatorTok{[}\StringTok{"execve"}\OperatorTok{]}
\DataTypeTok{comm\_in} \OperatorTok{=} \OperatorTok{[}\StringTok{"bash"}\OperatorTok{,} \StringTok{"sh"}\OperatorTok{,} \StringTok{"dash"}\OperatorTok{]}
\DataTypeTok{subject\_prefixes} \OperatorTok{=} \OperatorTok{[}\StringTok{"/tmp/"}\OperatorTok{]}
\end{Highlighting}
\end{Shaded}

\section{Detector Fields}\label{detector-fields}

Top-level fields:

\begin{itemize}
\tightlist
\item
  \texttt{id}: stable unique detector ID. Use lowercase words separated
  by \texttt{-}.
\item
  \texttt{name}: human-readable name shown in the UI.
\item
  \texttt{description}: short reason this detector exists.
\item
  \texttt{severity}: \texttt{info}, \texttt{warning}, or
  \texttt{critical}.
\item
  \texttt{tags}: list of strings for grouping and search.
\item
  \texttt{summary\_template}: detection summary. May reference event
  fields with \texttt{\{field\}}.
\item
  \texttt{threshold\_count}: optional count needed before emitting a
  detection.
\item
  \texttt{threshold\_window\_secs}: optional correlation window in
  seconds.
\item
  \texttt{group\_by}: optional event fields used as the threshold
  correlation key.
\end{itemize}

Match fields under \texttt{{[}match{]}}:

\begin{itemize}
\tightlist
\item
  \texttt{event\_types}: event names such as \texttt{execve},
  \texttt{file\_open}, \texttt{file\_write}, \texttt{ptrace},
  \texttt{module\_load}, \texttt{socket\_connect},
  \texttt{socket\_bind}, \texttt{setuid}, \texttt{mount}, \texttt{bpf}.
\item
  \texttt{actions}: \texttt{observe}, \texttt{allow}, or \texttt{block}.
\item
  \texttt{comm\_in}: exact process names from \texttt{comm}.
\item
  \texttt{comm\_prefixes}: process-name prefixes.
\item
  \texttt{subject\_prefixes}: path or subject prefixes.
\item
  \texttt{subject\_contains}: substrings in \texttt{subject}.
\item
  \texttt{dst\_ports}: destination or bind ports.
\item
  \texttt{dst\_ips}: destination IPv4 addresses.
\item
  \texttt{uids}: Linux UIDs.
\item
  \texttt{arg0\_in}: exact numeric \texttt{arg0} values.
\item
  \texttt{arg1\_in}: exact numeric \texttt{arg1} values.
\end{itemize}

All string matching is case-insensitive except \texttt{dst\_ips}.

\section{Event Fields}\label{event-fields}

Detectors can match or render these normalized fields:

\begin{itemize}
\tightlist
\item
  \texttt{agent\_id}
\item
  \texttt{ts}
\item
  \texttt{event\_type}
\item
  \texttt{action}
\item
  \texttt{severity}
\item
  \texttt{pid}
\item
  \texttt{uid}
\item
  \texttt{comm}
\item
  \texttt{dst\_ip}
\item
  \texttt{dst\_port}
\item
  \texttt{subject}
\item
  \texttt{arg0}
\item
  \texttt{arg1}
\item
  \texttt{ppid}
\item
  \texttt{ancestry}
\item
  \texttt{process\_path}
\item
  \texttt{pid\_ns}
\item
  \texttt{mount\_ns}
\item
  \texttt{net\_ns}
\end{itemize}

\section{Threshold Detectors}\label{threshold-detectors}

Use \texttt{threshold\_count}, \texttt{threshold\_window\_secs}, and
\texttt{group\_by} when one event is too noisy.

Example:

\begin{Shaded}
\begin{Highlighting}[]
\DataTypeTok{id} \OperatorTok{=} \StringTok{"repeated{-}bind{-}denials"}
\DataTypeTok{name} \OperatorTok{=} \StringTok{"Repeated Bind Denials"}
\DataTypeTok{description} \OperatorTok{=} \StringTok{"Flags repeated blocked bind attempts against the same port and process."}
\DataTypeTok{severity} \OperatorTok{=} \StringTok{"critical"}
\DataTypeTok{tags} \OperatorTok{=} \OperatorTok{[}\StringTok{"network"}\OperatorTok{,} \StringTok{"ips"}\OperatorTok{,} \StringTok{"threshold"}\OperatorTok{]}
\DataTypeTok{threshold\_count} \OperatorTok{=} \DecValTok{3}
\DataTypeTok{threshold\_window\_secs} \OperatorTok{=} \DecValTok{120}
\DataTypeTok{group\_by} \OperatorTok{=} \OperatorTok{[}\StringTok{"agent\_id"}\OperatorTok{,} \StringTok{"comm"}\OperatorTok{,} \StringTok{"dst\_port"}\OperatorTok{]}
\DataTypeTok{summary\_template} \OperatorTok{=} \StringTok{"\{comm\} hit \{match\_count\} blocked bind attempts on port \{dst\_port\} within 120s"}

\KeywordTok{[match]}
\DataTypeTok{event\_types} \OperatorTok{=} \OperatorTok{[}\StringTok{"socket\_bind"}\OperatorTok{]}
\DataTypeTok{actions} \OperatorTok{=} \OperatorTok{[}\StringTok{"block"}\OperatorTok{]}
\end{Highlighting}
\end{Shaded}

This emits on the third matching event in the 120 second window for the
same \texttt{agent\_id}, \texttt{comm}, and \texttt{dst\_port}.

\section{Block Detectors}\label{block-detectors}

LSM policy blocks are normal events with \texttt{action\ =\ "block"}.
Detector examples:

\begin{Shaded}
\begin{Highlighting}[]
\DataTypeTok{id} \OperatorTok{=} \StringTok{"blocked{-}exec{-}path"}
\DataTypeTok{name} \OperatorTok{=} \StringTok{"Blocked Exec Path"}
\DataTypeTok{description} \OperatorTok{=} \StringTok{"Highlights execution denied by selected exec path policy."}
\DataTypeTok{severity} \OperatorTok{=} \StringTok{"critical"}
\DataTypeTok{tags} \OperatorTok{=} \OperatorTok{[}\StringTok{"execution"}\OperatorTok{,} \StringTok{"ips"}\OperatorTok{,} \StringTok{"lsm"}\OperatorTok{]}
\DataTypeTok{summary\_template} \OperatorTok{=} \StringTok{"\{comm\} was blocked from executing \{subject\}"}

\KeywordTok{[match]}
\DataTypeTok{event\_types} \OperatorTok{=} \OperatorTok{[}\StringTok{"execve"}\OperatorTok{]}
\DataTypeTok{actions} \OperatorTok{=} \OperatorTok{[}\StringTok{"block"}\OperatorTok{]}
\end{Highlighting}
\end{Shaded}

\begin{Shaded}
\begin{Highlighting}[]
\DataTypeTok{id} \OperatorTok{=} \StringTok{"blocked{-}sensitive{-}write"}
\DataTypeTok{name} \OperatorTok{=} \StringTok{"Blocked Sensitive Write"}
\DataTypeTok{description} \OperatorTok{=} \StringTok{"Highlights denied writes to protected targets."}
\DataTypeTok{severity} \OperatorTok{=} \StringTok{"critical"}
\DataTypeTok{tags} \OperatorTok{=} \OperatorTok{[}\StringTok{"filesystem"}\OperatorTok{,} \StringTok{"ips"}\OperatorTok{,} \StringTok{"lsm"}\OperatorTok{]}
\DataTypeTok{summary\_template} \OperatorTok{=} \StringTok{"\{comm\} was blocked from writing target \{subject\}"}

\KeywordTok{[match]}
\DataTypeTok{event\_types} \OperatorTok{=} \OperatorTok{[}\StringTok{"file\_write"}\OperatorTok{]}
\DataTypeTok{actions} \OperatorTok{=} \OperatorTok{[}\StringTok{"block"}\OperatorTok{]}
\end{Highlighting}
\end{Shaded}

\section{Testing}\label{testing}

List loaded detectors:

\begin{Shaded}
\begin{Highlighting}[]
\ExtensionTok{curl} \AttributeTok{{-}s}\NormalTok{ http://127.0.0.1:8080/api/v1/detectors }\KeywordTok{|} \ExtensionTok{jq}
\end{Highlighting}
\end{Shaded}

List detections:

\begin{Shaded}
\begin{Highlighting}[]
\ExtensionTok{curl} \AttributeTok{{-}s}\NormalTok{ http://127.0.0.1:8080/api/v1/detections }\KeywordTok{|} \ExtensionTok{jq}
\end{Highlighting}
\end{Shaded}

Run replay tests after adding detector fixtures:

\begin{Shaded}
\begin{Highlighting}[]
\ExtensionTok{python3} \AttributeTok{{-}m}\NormalTok{ unittest tests.test\_replay}
\end{Highlighting}
\end{Shaded}

\section{Detector Limits}\label{detector-limits}

Current detectors are single-event or simple threshold rules. They do
not yet support:

\begin{itemize}
\tightlist
\item
  regex
\item
  negative matches
\item
  complex boolean logic
\item
  multi-stage sequence correlation
\item
  file hash matching
\item
  content scanning
\end{itemize}

Use agent or manager code changes for those features.

\clearpage

\chapter{Advanced Operations: Signature Antivirus with
LSM}\label{advanced-operations-signature-antivirus-with-lsm}

\begin{figure}[h]
\centering
\vspace{5cm}
\caption{Illustration mapping the logical boundary of the Advanced Operations: Signature Antivirus with LSM implementation.}
\end{figure}

A simple signature-based antivirus is possible in this project, but the
split matters:

\begin{itemize}
\tightlist
\item
  eBPF LSM is good at inline enforcement: allow or deny \texttt{exec},
  \texttt{open}, \texttt{write}, \texttt{ptrace}, and similar
  operations.
\item
  Userspace is the right place for signature scanning: hashing files,
  reading file contents, parsing signature databases, and making richer
  verdicts.
\end{itemize}

Do not try to build a full file-content scanner inside eBPF. BPF
programs are bounded, cannot freely read whole files, and are not a good
place for SHA/YARA-style work. Use BPF LSM to enforce verdicts that
userspace has already computed.

\section{Recommended Architecture}\label{recommended-architecture}

\begin{enumerate}
\def\labelenumi{\arabic{enumi}.}
\tightlist
\item
  Agent observes file activity and execution attempts.
\item
  Agent or helper scanner hashes files in userspace.
\item
  Agent checks hashes against a local signature database.
\item
  Agent writes malicious verdicts into BPF maps.
\item
  LSM hooks block future matching execution or file access.
\item
  Agent forwards \texttt{malware\_detected} and block events to the
  manager.
\item
  Manager stores detections and distributes signature updates.
\end{enumerate}

\section{Enforcement Keys}\label{enforcement-keys}

Prefer stable file identity over path-only matching:

\begin{itemize}
\tightlist
\item
  \texttt{dev}
\item
  \texttt{inode}
\item
  optional \texttt{ctime} or size
\item
  optional file hash
\end{itemize}

Path rules are useful for policy, but weak for antivirus because files
can move, link, or be replaced.

Good first implementation:

\begin{itemize}
\tightlist
\item
  Userspace scanner computes SHA-256.
\item
  Signature DB maps SHA-256 to malware name/severity.
\item
  Agent keeps a userspace cache:
  \texttt{(dev,\ inode)\ -\textgreater{}\ verdict}.
\item
  Agent syncs blocked \texttt{(dev,\ inode)} keys into an eBPF map.
\item
  \texttt{bprm\_check\_security} blocks execution if
  \texttt{(dev,\ inode)} is in the malicious map.
\item
  \texttt{file\_permission} can optionally block read or write for known
  malicious files.
\end{itemize}

\section{Signature Database Shape}\label{signature-database-shape}

Start simple:

\begin{Shaded}
\begin{Highlighting}[]
\FunctionTok{\{}
  \DataTypeTok{"version"}\FunctionTok{:} \DecValTok{1}\FunctionTok{,}
  \DataTypeTok{"updated\_at"}\FunctionTok{:} \StringTok{"2026{-}04{-}24T00:00:00Z"}\FunctionTok{,}
  \DataTypeTok{"hashes"}\FunctionTok{:} \OtherTok{[}
    \FunctionTok{\{}
      \DataTypeTok{"sha256"}\FunctionTok{:} \StringTok{"0123456789abcdef..."}\FunctionTok{,}
      \DataTypeTok{"name"}\FunctionTok{:} \StringTok{"EICAR{-}Test{-}File"}\FunctionTok{,}
      \DataTypeTok{"severity"}\FunctionTok{:} \StringTok{"critical"}\FunctionTok{,}
      \DataTypeTok{"action"}\FunctionTok{:} \StringTok{"block"}
    \FunctionTok{\}}
  \OtherTok{]}
\FunctionTok{\}}
\end{Highlighting}
\end{Shaded}

Manager can serve this as a policy endpoint later, for example:

\begin{Shaded}
\begin{Highlighting}[]
\NormalTok{GET /api/v1/signatures}
\end{Highlighting}
\end{Shaded}

Agent should cache the database locally so enforcement survives manager
outages.

\section{LSM Hooks To Use}\label{lsm-hooks-to-use}

Useful hooks:

\begin{itemize}
\tightlist
\item
  \texttt{bprm\_check\_security}: best first hook for blocking malicious
  executable launches.
\item
  \texttt{file\_permission}: can block reads/writes after a file is
  known malicious.
\item
  \texttt{inode\_create}, \texttt{inode\_rename},
  \texttt{inode\_unlink}, \texttt{file\_open}: possible future hooks for
  cache invalidation and quarantine workflows.
\end{itemize}

Best first milestone:

\begin{enumerate}
\def\labelenumi{\arabic{enumi}.}
\tightlist
\item
  Scan on \texttt{execve} observation or before
  \texttt{bprm\_check\_security} verdict is needed.
\item
  Cache clean/malicious file verdicts in userspace.
\item
  Block known-malicious \texttt{(dev,\ inode)} in
  \texttt{bprm\_check\_security}.
\end{enumerate}

\section{Race Conditions}\label{race-conditions}

Pure userspace scanning has a time-of-check/time-of-use risk:

\begin{itemize}
\tightlist
\item
  file is scanned clean
\item
  file content changes
\item
  file is executed before rescan
\end{itemize}

Reduce this with:

\begin{itemize}
\tightlist
\item
  cache key includes \texttt{mtime}, \texttt{ctime}, and size
\item
  invalidate cache on write events
\item
  rescan on first execution after write
\item
  block unknown execution until scan completes if strict mode is enabled
\end{itemize}

Strict mode is safer but can break usability if scanners are slow.

\section{Scanner Options}\label{scanner-options}

Minimal scanner:

\begin{itemize}
\tightlist
\item
  SHA-256 exact hash signatures
\item
  local JSON signature DB
\item
  verdict cache
\item
  BPF map sync for malicious file IDs
\end{itemize}

Better scanner:

\begin{itemize}
\tightlist
\item
  SHA-256 plus fuzzy hashes
\item
  YARA-compatible rules in userspace
\item
  quarantine metadata
\item
  allowlists for trusted paths/signers
\item
  per-agent signature policy from manager
\end{itemize}

\section{Event Types To Add}\label{event-types-to-add}

Suggested manager event types:

\begin{itemize}
\tightlist
\item
  \texttt{malware\_scan}: scan completed
\item
  \texttt{malware\_detected}: malicious signature matched
\item
  \texttt{malware\_block}: LSM blocked access to known malicious file
\end{itemize}

Useful event fields:

\begin{itemize}
\tightlist
\item
  \texttt{subject}: path if available
\item
  \texttt{arg0}: signature database version or scanner verdict code
\item
  \texttt{arg1}: file size or signature rule ID if numeric
\item
  \texttt{process\_path}: process that touched the file
\item
  additional raw JSON fields: \texttt{sha256}, \texttt{signature\_name},
  \texttt{file\_dev}, \texttt{file\_inode}
\end{itemize}

\section{Practical Next Step}\label{practical-next-step}

Build exact-hash blocking first:

\begin{enumerate}
\def\labelenumi{\arabic{enumi}.}
\tightlist
\item
  Add a manager signature DB endpoint or local agent signature file.
\item
  Add agent userspace SHA-256 scanner and verdict cache.
\item
  Add BPF map keyed by \texttt{(dev,\ inode)} for malicious files.
\item
  Add \texttt{bprm\_check\_security} enforcement against that map.
\item
  Emit \texttt{malware\_detected} and \texttt{malware\_block} events.
\item
  Add TOML detectors for those event types.
\end{enumerate}

This gives antivirus-like blocking while keeping BPF small and reliable.

\clearpage

\chapter{Future Roadmap \& Extension
Plan}\label{future-roadmap-extension-plan}

\begin{figure}[h]
\centering
\vspace{5cm}
\caption{Illustration mapping the logical boundary of the Future Roadmap & Extension Plan implementation.}
\end{figure}

This is the next-phase plan for KShield after the detector engine and
expanded LSM policy work now in the repo.

\section{Phase 1: Completed in This
Iteration}\label{phase-1-completed-in-this-iteration}

\begin{itemize}
\tightlist
\item
  Add a universal detector definition format using TOML.
\item
  Persist detector hits separately from raw events.
\item
  Expand telemetry beyond \texttt{execve} into file access, ptrace, and
  module load attempts.
\item
  Expand LSM enforcement beyond \texttt{socket\_connect} into
  \texttt{socket\_bind}.
\item
  Surface detector hits and detector inventory in the dashboard.
\end{itemize}

\section{Phase 2: Kernel Telemetry
Expansion}\label{phase-2-kernel-telemetry-expansion}

\begin{itemize}
\tightlist
\item
  Add process ancestry and parent PID to event payloads.
\item
  Add \texttt{setuid}, \texttt{mount}, and \texttt{bpf} syscall
  visibility.
\item
  Capture richer file metadata for sensitive paths.
\item
  Add optional container and namespace context when available.
\end{itemize}

\section{Phase 3: Detector Semantics}\label{phase-3-detector-semantics}

\begin{itemize}
\tightlist
\item
  Support multi-event correlation windows, not just single-event
  matches.
\item
  Add threshold detectors such as repeated failed binds or bursty ptrace
  attempts.
\item
  Add allowlists and suppression rules per detector.
\item
  Add detector test fixtures and offline replay of saved event streams.
\end{itemize}

\section{Phase 4: Enforcement}\label{phase-4-enforcement}

\begin{itemize}
\tightlist
\item
  Add executable deny rules for selected paths or hashes.
\item
  Add sensitive-file write protection.
\item
  Add optional ptrace deny rules.
\item
  Add scoped policy application per agent rather than global policy
  only.
\end{itemize}

\section{Phase 5: Operations and UX}\label{phase-5-operations-and-ux}

\begin{itemize}
\tightlist
\item
  Turn the manager into a persistent systemd service on the host.
\item
  Add dashboard filters for detector ID, severity, and agent.
\item
  Add event-to-detection linking in the UI.
\item
  Add export endpoints for JSON and CSV.
\item
  Add health and lag indicators for detector reloads and policy syncs.
\end{itemize}

\clearpage

\end{document}
